\begin{textsample}{POS Dim 1 – 2003 – Score 24.00 – 2003/t000436.txt}  \label{ex:f1_pos_002}
I ' m supposed to scare you, because it ' s about fear, right?
And you should be really afraid, but not for \textbf{the} reasons why you think you should be.
You should be really afraid that — if we stick up \textbf{the} first \textbf{slide} on this thing — there we go — that you ' re missing out.
Because if you spend this week thinking about Iraq and thinking about Bush and thinking about \textbf{the} stock market, you ' re going to miss one of \textbf{the} greatest adventures that we ' ve ever been on.
And this is what this adventure ' s really about.
This is crystallized DNA.
Every life form on this \textbf{planet} — every insect, every bacteria, every plant, every animal, every human, every politician — is coded in that stuff.
And if you want to take a single crystal of DNA, it looks like that.
And we ' re just beginning to understand this stuff.
And this is \textbf{the} single most exciting adventure that we have ever been on.
It ' s \textbf{the} single greatest mapping project we ' ve ever been on.
If you think that \textbf{the} mapping of America ' s made a difference, or landing on \textbf{the} moon, or this other stuff, it ' s \textbf{the} map of ourselves and \textbf{the} map of every plant and every insect and every bacteria that really makes a difference.
And it ' s beginning to tell us a lot about \textbf{evolution}.
It turns out that what this stuff is — and Richard Dawkins has written about this — is, this is really a river out of Eden.
So, \textbf{the} 3. 2 billion base pairs inside each of your cells is really a history of where you ' ve been for \textbf{the} past billion years.
And we could start dating things, and we could start changing medicine and archeology.
It turns out that if you take \textbf{the} human species about 700 years ago, white Europeans diverged from black Africans in a very significant way.
White Europeans were subject to \textbf{the} plague.
And when they were subject to \textbf{the} plague, most people didn ' t survive, but those who survived had a mutation on \textbf{the} CCR5 receptor.
And that mutation was passed on to their kids because they ' re \textbf{the} ones that survived, so there was a great deal of population pressure.
In Africa, because you didn ' t have these cities, you didn ' t have that CCR5 population pressure mutation.
We can date it to 700 years ago.
That is one of \textbf{the} reasons why AIDS is raging across Africa as fast as it is, and not as fast across Europe.
And we ' re beginning to find these little things for \textbf{malaria}, for sickle cell, for cancers.
And in \textbf{the} measure that we map ourselves, this is \textbf{the} single greatest adventure that we ' ll ever be on.
And this Friday, I want you to pull out a really good bottle of wine, and I want you to toast these two people.
Because this Friday, 50 years ago, Watson and Crick found \textbf{the} structure of DNA, and that is almost as important a date as \textbf{the} 12th of February when we first mapped ourselves, but anyway, we ' ll get to that.
I thought we ' d talk about \textbf{the} new zoo.
So, all you guys have heard about DNA, all \textbf{the} stuff that DNA does, but some of \textbf{the} stuff we ' re discovering is kind of nifty because this turns out to be \textbf{the} single most abundant species on \textbf{the} \textbf{planet}.
If you think you ' re successful or \textbf{cockroaches} are successful, it turns out that there ' s ten trillion trillion Pleurococcus sitting out there.
And we didn ' t know that Pleurococcus was out there, which is part of \textbf{the} reason why this whole species-mapping project is so important.
Because we ' re just beginning to learn where we came from and what we are.
And we ' re finding amoebas like this.
This is \textbf{the} amoeba dubia.
And \textbf{the} amoeba dubia doesn ' t look like much, except that each of you has about 3. 2 billion letters, which is what makes you you, as far as gene code inside each of your cells, and this little amoeba which, you know, sits in water in hundreds and millions and billions, turns out to have 620 billion base pairs of gene code inside.
So, this little thingamajig has a \textbf{genome} that ' s 200 times \textbf{the} size of yours.
And if you ' re thinking of efficient information storage mechanisms, it may not turn out to be chips.
It may turn out to be something that looks a little like that amoeba.
And, again, we ' re learning from life and how life works.
This funky little thing : people didn ' t used to think that it was worth taking samples out of nuclear reactors because it was dangerous and, of course, nothing lived there.
And then finally somebody picked up a microscope and looked at \textbf{the} water that was sitting next to \textbf{the} cores.
And sitting next to that water in \textbf{the} cores was this little Deinococcus radiodurans, doing a backstroke, having its chromosomes blown apart every day, six, seven times, restitching them, living in about 200 times \textbf{the} radiation that would kill you.
And by now you should be getting a hint as to how diverse and how important and how interesting this journey into life is, and how many different life forms there are, and how there can be different life forms living in very different places, maybe even outside of this \textbf{planet}.
Because if you can live in radiation that looks like this, that brings up a whole series of interesting questions.
This little thingamajig : we didn ' t know this thingamajig existed.
We should have known that this existed because this is \textbf{the} only bacteria that you can see to \textbf{the} naked eye.
So, this thing is 0. 75 millimeters.
It lives in a deep trench off \textbf{the} coast of Namibia.
And what you ' re looking at with this namibiensis is \textbf{the} biggest bacteria we ' ve ever seen.
So, it ' s about \textbf{the} size of a little period on a sentence.
Again, we didn ' t know this thing was there three years ago.
We ' re just beginning this journey of life in \textbf{the} new zoo.
This is a really odd one.
This is Ferroplasma. \textbf{The} reason why Ferroplasma is interesting is because it eats \textbf{iron}, lives inside \textbf{the} equivalent of battery acid, and excretes sulfuric acid.
So, when you think of odd life forms, when you think of what it takes to live, it turns out this is a very efficient life form, and they call it an archaea.
Archaea means"\textbf{the} ancient ones."And \textbf{the} reason why they ' re ancient is because this thing came up when this \textbf{planet} was covered by things like sulfuric acid in batteries, and it was eating \textbf{iron} when \textbf{the} earth was part of a melted core.
So, it ' s not just dogs and cats and \textbf{whales} and dolphins that you should be aware of and interested in on this little journey.
Your fear should be that you are not, that you ' re paying attention to stuff which is temporal.
I mean, George Bush — he ' s going to be gone, alright?
Life isn ' t.
Whether \textbf{the} humans survive or don ' t survive, these things are going to be living on this \textbf{planet} or other \textbf{planets}.
And it ' s just beginning to understand this code of DNA that ' s really \textbf{the} most exciting intellectual adventure that we ' ve ever been on.
And you can do strange things with this stuff.
This is a baby gaur.
Conservation group gets together, tries to figure out how to breed an animal that ' s almost extinct.
They can ' t do it naturally, so what they do with this thing is they take a spoon, take some cells out of an adult gaur ' s mouth, code, take \textbf{the} cells from that and insert it into a fertilized cow ' s egg, reprogram cow ' s egg — different gene code.
When you do that, \textbf{the} cow gives birth to a gaur.
We are now experimenting with bongos, pandas, elands, Sumatran \textbf{tigers}, and \textbf{the} Australians — bless their hearts — are playing with these things.
Now, \textbf{the} last of these things died in September 1936.
These are Tasmanian \textbf{tigers}. \textbf{The} last known one died at \textbf{the} Hobart Zoo.
But it turns out that as we learn more about gene code and how to reprogram species, we may be able to close \textbf{the} gene gaps in deteriorate DNA.
And when we learn how to close \textbf{the} gene gaps, then we can put a full string of DNA together.
And if we do that, and insert this into a fertilized wolf ' s egg, we may give birth to an animal that hasn ' t walked \textbf{the} earth since 1936.
And then you can start going back further, and you can start thinking about dodos, and you can think about other species.
And in other places, like Maryland, they ' re trying to figure out what \textbf{the} primordial ancestor is.
Because each of us contains our entire gene code of where we ' ve been for \textbf{the} past billion years, because we ' ve \textbf{evolved} from that stuff, you can take that tree of life and collapse it back, and in \textbf{the} measure that you learn to reprogram, maybe we ' ll give birth to something that is very close to \textbf{the} first primordial ooze.
And it ' s all coming out of things that look like this.
These are companies that didn ' t exist five years ago.
Huge gene sequencing facilities \textbf{the} size of football fields.
Some are public.
Some are private.
It takes about 5 billion dollars to sequence a human being \textbf{the} first time.
Takes about 3 million dollars \textbf{the} second time.
We will have a 1, 000-dollar \textbf{genome} within \textbf{the} next five to eight years.
That means each of you will contain on a CD your entire gene code.
And it will be really boring.
It will read like this. \textbf{The} really neat thing about this stuff is that ' s life.
And Laurie ' s going to talk about this one a little bit.
Because if you happen to find this one inside your body, you ' re in big trouble, because that ' s \textbf{the} source code for Ebola.
That ' s one of \textbf{the} deadliest diseases known to humans.
But plants work \textbf{the} same way and insects work \textbf{the} same way, and this apple works \textbf{the} same way.
This apple is \textbf{the} same thing as this floppy \textbf{disk}.
Because this thing codes ones and zeros, and this thing codes A, T, C, Gs, and it sits up there, absorbing energy on a tree, and one fine day it has enough energy to say, execute, and it goes [ thump ].
Right?
And when it does that, pushes a.
EXE, what it does is, it executes \textbf{the} first line of code, which reads just like that, AATCAGGGACCC, and that means : make a root.
Next line of code : make a stem.
Next line of code, TACGGGG : make a flower that ' s white, that blooms in \textbf{the} \textbf{spring}, that smells like this.
In \textbf{the} measure that you have \textbf{the} code and \textbf{the} measure that you read it — and, by \textbf{the} way, \textbf{the} first plant was read two years ago ; \textbf{the} first human was read two years ago ; \textbf{the} first insect was read two years ago. \textbf{The} first thing that we ever read was in 1995 : a little bacteria called Haemophilus influenzae.
In \textbf{the} measure that you have \textbf{the} source code, as all of you know, you can change \textbf{the} source code, and you can reprogram life forms so that this little thingy becomes a vaccine, or this little thingy starts producing biomaterials, which is why DuPont is now growing a form of polyester that feels like silk in corn.
This changes all rules.
This is life, but we ' re reprogramming it.
This is what you look like.
This is one of your chromosomes.
And what you can do now is, you can outlay exactly what your chromosome is, and what \textbf{the} gene code on that chromosome is right here, and what those genes code for, and what animals they code against, and then you can tie it to \textbf{the} literature.
And in \textbf{the} measure that you can do that, you can go home today, and get on \textbf{the} Internet, and access \textbf{the} world ' s biggest public library, which is a library of life.
And you can do some pretty strange things because in \textbf{the} same way as you can reprogram this apple, if you go to Cliff Tabin ' s lab at \textbf{the} Harvard Medical School, he ' s reprogramming chicken embryos to grow more \textbf{wings}.
Why would Cliff be doing that?
He doesn ' t have a restaurant. \textbf{The} reason why he ' s reprogramming that animal to have more \textbf{wings} is because when you used to play with lizards as a little child, and you picked up \textbf{the} lizard, sometimes \textbf{the} tail fell off, but it regrew.
Not so in human beings : you cut off an arm, you cut off a leg — it doesn ' t regrow.
But because each of your cells contains your entire gene code, each cell can be reprogrammed, if we don ' t stop stem cell research and if we don ' t stop genomic research, to express different body functions.
And in \textbf{the} measure that we learn how chickens grow \textbf{wings}, and what \textbf{the} program is for those cells to differentiate, one of \textbf{the} things we ' re going to be able to do is to stop undifferentiated cells, which you know as cancer, and one of \textbf{the} things we ' re going to learn how to do is how to reprogram cells like stem cells in such a way that they express bone, stomach, skin, pancreas.
And you are likely to be wandering around — and your children — on regrown body parts in a reasonable period of time, in some places in \textbf{the} world where they don ' t stop \textbf{the} research.
How ' s this stuff work?
If each of you differs from \textbf{the} person next to you by one in a thousand, but only three percent codes, which means it ' s only one in a thousand times three percent, very small differences in expression and punctuation can make a significant difference.
Take a simple declarative sentence.
Right?
That ' s perfectly clear.
So, men read that sentence, and they look at that sentence, and they read this. \textbf{Okay}?
Now, women look at that sentence and they say, uh-uh, wrong.
This is \textbf{the} way it should be seen.
That ' s what your genes are doing.
That ' s why you differ from this person over here by one in a thousand.
Right?
But, you know, he ' s reasonably good looking, but...
I won ' t go there.
You can do this stuff even without changing \textbf{the} punctuation.
You can look at this, right?
And they look at \textbf{the} world a little differently.
They look at \textbf{the} same world and they say...
That ' s how \textbf{the} same gene code — that ' s why you have 30, 000 genes, mice have 30, 000 genes, husbands have 30, 000 genes.
Mice and men are \textbf{the} same.
Wives know that, but anyway.
You can make very small changes in gene code and get really different outcomes, even with \textbf{the} same string of letters.
That ' s what your genes are doing every day.
That ' s why sometimes a person ' s genes don ' t have to change a lot to get cancer.
These little chippies, these things are \textbf{the} size of a credit card.
They will test any one of you for 60, 000 genetic conditions.
That brings up questions of privacy and insurability and all kinds of stuff, but it also allows us to start going after diseases, because if you run a person who has leukemia through something like this, it turns out that three diseases with completely similar clinical syndromes are completely different diseases.
Because in ALL leukemia, that set of genes over there over-expresses.
In MLL, it ' s \textbf{the} middle set of genes, and in AML, it ' s \textbf{the} bottom set of genes.
And if one of those particular things is expressing in your body, then you take Gleevec and you ' re cured.
If it is not expressing in your body, if you don ' t have one of those types — a particular one of those types — don ' t take Gleevec.
It won ' t do anything for you.
Same thing with Receptin if you ' ve got breast cancer.
Don ' t have an HER- 2 receptor?
Don ' t take Receptin.
Changes \textbf{the} nature of medicine.
Changes \textbf{the} predictions of medicine.
Changes \textbf{the} way medicine works. \textbf{The} greatest repository of knowledge when most of us went to college was this thing, and it turns out that this is not so important any more. \textbf{The} U.
S.
Library of Congress, in terms of its printed volume of data, contains less data than is coming out of a good genomics company every month on a compound basis.
Let me say that again : A single genomics company generates more data in a month, on a compound basis, than is in \textbf{the} printed \textbf{collections} of \textbf{the} Library of Congress.
This is what ' s been powering \textbf{the} U.
S. economy.
It ' s Moore ' s Law.
So, all of you know that \textbf{the} price of \textbf{computers} halves every 18 months and \textbf{the} power doubles, right?
Except that when you lay that side by side with \textbf{the} speed with which gene data ' s being deposited in GenBank, Moore ' s Law is right here : it ' s \textbf{the} \textbf{blue} line.
This is on a log scale, and that ' s what superexponential growth means.
This is going to push \textbf{computers} to have to grow faster than they ' ve been growing, because so far, there haven ' t been applications that have been required that need to go faster than Moore ' s Law.
This stuff does.
And here ' s an interesting map.
This is a map which was finished at \textbf{the} Harvard Business School.
One of \textbf{the} really interesting questions is, if all this data ' s free, who ' s using it?
This is \textbf{the} greatest public library in \textbf{the} world.
Well, it turns out that there ' s about 27 trillion bits moving inside from \textbf{the} United States to \textbf{the} United States ; about 4. 6 trillion is going over to those European countries ; about 5. 5 ' s going to Japan ; there ' s almost no communication between Japan, and nobody else is literate in this stuff.
It ' s free.
No one ' s reading it.
They ' re focusing on \textbf{the} war ; they ' re focusing on Bush ; they ' re not interested in life.
So, this is what a new map of \textbf{the} world looks like.
That is \textbf{the} genomically literate world.
And that is a problem.
In fact, it ' s not a genomically literate world.
You can break this out by states.
And you can watch states rise and fall depending on their ability to speak a language of life, and you can watch New York fall off a cliff, and you can watch New Jersey fall off a cliff, and you can watch \textbf{the} rise of \textbf{the} new empires of intelligence.
And you can break it out by counties, because it ' s specific counties.
And if you want to get more specific, it ' s actually specific zip codes.
So, you want to know where life is happening?
Well, in Southern California it ' s happening in 92121.
And that ' s it.
And that ' s \textbf{the} \textbf{triangle} between Salk, Scripps, UCSD, and it ' s called Torrey Pines Road.
That means you don ' t need to be a big nation to be successful ; it means you don ' t need a lot of people to be successful ; and it means you can move most of \textbf{the} wealth of a country in about three or four carefully picked 747s.
Same thing in Massachusetts.
Looks more spread out but — oh, by \textbf{the} way, \textbf{the} ones that are \textbf{the} same color are contiguous.
What ' s \textbf{the} net effect of this?
In an agricultural society, \textbf{the} difference between \textbf{the} richest and \textbf{the} poorest, \textbf{the} most productive and \textbf{the} least productive, was five to one.
Why?
Because in \textbf{agriculture}, if you had 10 kids and you grow up a little bit earlier and you work a little bit harder, you could produce about five times more wealth, on average, than your neighbor.
In a knowledge society, that number is now 427 to 1.
It really matters if you ' re literate, not just in reading and writing in English and French and German, but in Microsoft and Linux and Apple.
And very soon it ' s going to matter if you ' re literate in life code.
So, if there is something you should fear, it ' s that you ' re not keeping your eye on \textbf{the} ball.
Because it really matters who speaks life.
That ' s why nations rise and fall.
And it turns out that if you went back to \textbf{the} 1870s, \textbf{the} most productive nation on earth was Australia, per person.
And New Zealand was way up there.
And then \textbf{the} U.
S. came in about 1950, and then Switzerland about 1973, and then \textbf{the} U.
S. got back on top — beat up their chocolates and cuckoo clocks.
And today, of course, you all know that \textbf{the} most productive nation on earth is Luxembourg, producing about one third more wealth per person per year than America.
Tiny landlocked state.
No \textbf{oil}.
No diamonds.
No natural resources.
Just smart people moving bits.
Different rules.
Here ' s differential productivity rates.
Here ' s how many people it takes to produce a single U.
S. \textbf{patent}.
So, about 3, 000 Americans, 6, 000 Koreans, 14, 000 Brits, 790, 000 Argentines.
You want to know why Argentina ' s crashing?
It ' s got nothing to do with inflation.
It ' s got nothing to do with privatization.
You can take a Harvard-educated Ivy League economist, stick him in charge of Argentina.
He still crashes \textbf{the} country because he doesn ' t understand how \textbf{the} rules have changed.
Oh, yeah, and it takes about 5. 6 million Indians.
Well, watch what happens to India.
India and China used to be 40 percent of \textbf{the} global economy just at \textbf{the} Industrial Revolution, and they are now about 4. 8 percent.
Two billion people.
One third of \textbf{the} global population producing 5 percent of \textbf{the} wealth because they didn ' t get this change, because they kept treating their people like serfs instead of like shareholders of a common project.
They didn ' t keep \textbf{the} people who were educated.
They didn ' t foment \textbf{the} businesses.
They didn ' t do \textbf{the} IPOs.
Silicon Valley did.
And that ' s why they say that Silicon Valley has been powered by ICs.
Not integrated circuits : Indians and Chinese.
Here ' s what ' s happening in \textbf{the} world.
It turns out that if you ' d gone to \textbf{the} U.
N. in 1950, when it was founded, there were 50 countries in this world.
It turns out there ' s now about 192.
Country after country is splitting, seceding, succeeding, failing — and it ' s all getting very fragmented.
And this has not stopped.
In \textbf{the} 1990s, these are sovereign states that did not exist before 1990.
And this doesn ' t include \textbf{fusions} or name changes or changes in flags.
We ' re generating about 3. 12 states per year.
People are taking control of their own states, sometimes for \textbf{the} better and sometimes for \textbf{the} worse.
And \textbf{the} really interesting thing is, you and your kids are empowered to build great empires, and you don ' t need a lot to do it.
And, given that \textbf{the} music is over, I was going to talk about how you can use this to generate a lot of wealth, and how code works.
Moderator : Two minutes.
Juan Enriquez : No, I ' m going to stop there and we ' ll do it next year because I don ' t want to take any of Laurie ' s time.
But thank you very much.

% matched lemmas: agriculture, blue, cockroach, collection, computer, disk, evolution, evolve, fusion, genome, iron, malaria, oil, okay, patent, planet, slide, spring, the, tiger, triangle, whale, wing
\end{textsample}
